% Chapter fragment: load amsmath and graphicx in the main report.
% Derivation follows ../MPC/transferFunction4TankSystem.md and the
% implementation in ../MPC/four-tank-model/problem3*.m, checked against common/.
\chapter{PID Control}
\label{chap:problem3}

\section{Linear Model and Input--Output Pairing}

The purpose of feedback is to regulate the two lower-tank levels using
pump flows $F_1$ and $F_2$. The nonlinear mass balances from
Chapter~\ref{chap:problem2} are linearized about a positive equilibrium
$(x_s,u_s,d_s)$ satisfying $f(x_s,u_s,d_s,p)=0$. Disturbances are held at
their nominal values during controller design. These values need not be
zero: the experiments below use different operating points.

Following the analytical derivation in
\texttt{../MPC/transferFunction4TankSystem.md}, the equilibrium flows,
levels, and masses satisfy
\begin{align}
 q_{3,s}&=(1-\gamma_2)F_{2,s}+F_{3,s},&
 q_{4,s}&=(1-\gamma_1)F_{1,s}+F_{4,s},\\
 q_{1,s}&=\gamma_1F_{1,s}+q_{3,s},&
 q_{2,s}&=\gamma_2F_{2,s}+q_{4,s},\\
 h_{i,s}&=\frac{(q_{i,s}/a_i)^2}{2g},&
 m_{i,s}&=\rho A_i h_{i,s}.
\end{align}
The tangent-line approximation of each outlet flow is
\begin{equation}
 q_i(m_i)\simeq q_{i,s}+k_i\delta m_i,\qquad
 k_i=\left.\frac{dq_i}{dm_i}\right|_{m_{i,s}}
     =\frac{q_{i,s}}{2m_{i,s}},\qquad
 T_i=\frac{1}{\rho k_i}=\frac{2A_i h_{i,s}}{q_{i,s}}.
\end{equation}
With $\delta x=x-x_s$, $\delta u=u-u_s$, and $\delta d=d-d_s$,
\begin{equation}
 \delta\dot x=A\delta x+B\delta u+E\delta d,\qquad
 \delta z=C_h\delta x,
\end{equation}
where
\begin{equation}
 A=\begin{bmatrix}
 -1/T_1&0&1/T_3&0\\0&-1/T_2&0&1/T_4\\
 0&0&-1/T_3&0\\0&0&0&-1/T_4
 \end{bmatrix},\quad B=\rho U,\quad E=\rho D,
 \quad C_h=\begin{bmatrix}(\rho A_1)^{-1}&0&0&0\\
 0&(\rho A_2)^{-1}&0&0\end{bmatrix}.
\end{equation}
The routing matrices $U,D$ are those of Chapter~\ref{chap:problem2}.
All plant poles are negative at a positive equilibrium, but the
linear approximation remains local because discharge gains vary with level.

For zero initial deviations and $\delta d=0$, the upper-tank equations give
\begin{equation}
 \delta m_3(s)=\frac{\rho(1-\gamma_2)T_3}{1+sT_3}\delta F_2(s),\qquad
 \delta m_4(s)=\frac{\rho(1-\gamma_1)T_4}{1+sT_4}\delta F_1(s).
\end{equation}
Substituting into the lower-tank equations yields
\begin{equation}
 G_m(s)=\begin{bmatrix}
 \dfrac{\rho\gamma_1T_1}{1+sT_1}&
 \dfrac{\rho(1-\gamma_2)T_1}{(1+sT_1)(1+sT_3)}\\[6pt]
 \dfrac{\rho(1-\gamma_1)T_2}{(1+sT_2)(1+sT_4)}&
 \dfrac{\rho\gamma_2T_2}{1+sT_2}
 \end{bmatrix},\qquad
 G_h(s)=\operatorname{diag}\left(\frac{1}{\rho A_1},
 \frac{1}{\rho A_2}\right)G_m(s).
 \label{eq:p3-transfer}
\end{equation}
The direct paths $F_1\to h_1$ and $F_2\to h_2$ have one tank lag;
the cross paths have two lags. This supports pairing pump 1 with tank 1
and pump 2 with tank 2, although the loops interact.

A quantitative check uses the relative gain array (RGA), with
$\circ$ denoting elementwise multiplication:
\begin{equation}
 \Lambda=G_h(0)\circ\bigl(G_h(0)^{-1}\bigr)^T
 =\frac{1}{\gamma_1+\gamma_2-1}
 \begin{bmatrix}
 \gamma_1\gamma_2&-(1-\gamma_1)(1-\gamma_2)\\
 -(1-\gamma_1)(1-\gamma_2)&\gamma_1\gamma_2
 \end{bmatrix}
 =\begin{bmatrix}1.4&-0.4\\-0.4&1.4\end{bmatrix}.
\end{equation}
The positive diagonal relative gains favour the direct pairing; negative
cross gains make the alternative pairing unattractive. The diagonal value
1.4 also indicates appreciable steady-state interaction.

\section{Static Decoupling and Controller Equations}

Let $G$ denote either $G_m$ for mass feedback or $G_h$ for level feedback.
The static decoupler and effective plant are
\begin{equation}
 W=G(0)^{-1},\qquad G_{\mathrm{eff}}(s)=G(s)W,\qquad
 G_{\mathrm{eff}}(0)=I_2.
 \label{eq:p3-decoupler}
\end{equation}
The inverse exists because $\gamma_1+\gamma_2-1=0.3\ne0$.
Numerically, the implementation solves $G(0)W=I_2$ using MATLAB's
backslash operator. Static decoupling removes interaction at zero
frequency; $G_{\mathrm{eff}}(s)$ generally retains off-diagonal dynamics.
It therefore does not produce two fully independent SISO plants.

For error $e=r-y$ and virtual control $v$, the control law is
\begin{equation}
 v(s)=C_{\mathrm{diag}}(s)e(s),\qquad
 u=u_s+Wv,\qquad C_{\mathrm{diag}}=\operatorname{diag}(C_1,C_2).
\end{equation}
The standard controller structures are
\begin{equation}
 C_i^{\mathrm P}(s)=K_{p,i},\qquad
 C_i^{\mathrm{PI}}(s)=K_{p,i}+\frac{K_{i,i}}{s},\qquad
 C_i^{\mathrm{PID}}(s)=K_{p,i}+\frac{K_{i,i}}{s}
                       +\frac{K_{d,i}s}{1+T_{f,i}s}.
\end{equation}
The derivative filter $T_{f,i}>0$ limits high-frequency noise amplification.
For an ideal unit-static-gain effective plant, P control leaves the
constant-reference error $e_{i,\infty}=r_i/(1+K_{p,i})$ in deviation
coordinates. Integral action can remove constant offsets when the closed
loop is stable and the requested equilibrium is feasible. Derivative
feedback may improve transients, but its noise sensitivity must be evaluated.

The implemented experiments use PI controllers tuned independently on
$G_{\mathrm{eff},11}$ and $G_{\mathrm{eff},22}$ with MATLAB's
\texttt{pidtune} and \texttt{DesignFocus='reference-tracking'}.
The complete multivariable linear closed loop is checked using
\begin{equation}
 K=WC_{\mathrm{diag}},\qquad L=GK,\qquad
 T_{\mathrm{cl}}=(I_2+L)^{-1}L.
\end{equation}
Its poles are stable for the three designs reported here, and its static
reference gain is $I_2$ to numerical precision. Tuning the diagonal channels
alone does not guarantee stability of the interacting multivariable loop.

\begin{table}[htbp]
 \centering
 \caption{Operating points and PI gains. Flows are in $\mathrm{cm^3/s}$.
 Gains act on virtual channels normalized to unit static gain; $K_p$ is
 dimensionless and $K_i$ is in $\mathrm{s^{-1}}$.}
 \label{tab:p3-gains}
 \begin{tabular}{lrrrrrr}
 \hline
 Experiment & $u_s^T$ & $d_s^T$ & $K_{p,1}$ & $K_{i,1}$ & $K_{p,2}$ & $K_{i,2}$\\
 \hline
 Analysis & $(100,300)$ & $(100,100)$ & 1.05838 & 0.068357 & 1.03345 & 0.052202\\
 Level control & $(300,300)$ & $(100,100)$ & 0.99280 & 0.046257 & 0.99597 & 0.039951\\
 Comparison & $(300,300)$ & $(0,0)$ & 0.98203 & 0.063979 & 1.00857 & 0.054040\\
 \hline
 \end{tabular}
\end{table}
The analysis and comparison use mass-coordinate decouplers; the level
experiment uses a level-coordinate decoupler. These matrices have different
units and must be used with the corresponding error coordinates. In the
level experiment,
\begin{equation}
 W_h\simeq\begin{bmatrix}9.3170&-3.4358\\-5.3240&6.8717\end{bmatrix}
 \quad [\mathrm{cm^3/(s\,cm)}].
\end{equation}

\section{Closed-Loop Experiments and Results}

\subsection{Linear step responses}

The analysis operating point has levels
$h_s=[21.152,41.458,12.217,6.633]^T\ \mathrm{cm}$.
Figure~\ref{fig:p3-open} shows linear unit pump-step responses in mass
coordinates. Direct responses are first order and cross responses are
slower cascades, as predicted by \eqref{eq:p3-transfer}.
Figure~\ref{fig:p3-linear} shows the full PI-controlled linear reference
response. Off-diagonal responses illustrate the transient coupling that
remains after static decoupling.

\subsection{Nonlinear level regulation}

The nonlinear level experiment starts at equilibrium for
$u_s=[300,300]^T$ and $d_s=[100,100]^T$, with lower-tank levels
$[46.331,62.576]^T\ \mathrm{cm}$. The reference is set to
$[45,30]^T\ \mathrm{cm}$ at the beginning of the simulation. The continuous
PI controllers are discretized by Tustin's substitution,
\begin{equation}
 s\simeq\frac{2}{T_s}\frac{z-1}{z+1},\qquad
 C_i(z)=K_{p,i}+\frac{K_{i,i}T_s}{2}\frac{z+1}{z-1},\qquad T_s=2\ \mathrm{s}.
\end{equation}
MATLAB state-space realizations of these discrete controllers update the
controller states. Pump commands are held between samples, clipped below
at zero, and the nonlinear plant is integrated with \texttt{ode45}.
There is no upper pump bound or anti-windup mechanism in this experiment.
Measurement noise is zero and disturbances remain constant.

Figure~\ref{fig:p3-level} shows the levels and pump commands. At
$1200\ \mathrm{s}$ the true levels are
$[44.999284,29.999880]^T\ \mathrm{cm}$ and the pump flows are
$[419.766,42.587]^T\ \mathrm{cm^3/s}$.
The residual errors are below $0.001\ \mathrm{cm}$ in both tanks.
Using a settling band of 2\% of each initial reference change, the
sampled settling times are $774\ \mathrm{s}$ for tank 1 and
$444\ \mathrm{s}$ for tank 2. Tank 1 first rises despite its lower
reference, and tank 2 undershoots before settling; these transients show
the effects of residual interaction and actuator saturation.
During the transient, pump 2 reaches its zero lower bound. The recorded
pump ranges are $385.260\leq F_1\leq435.608$ and
$0\leq F_2\leq75.472\ \mathrm{cm^3/s}$.
Although the target in tank 1 changes only slightly, its pump must increase
substantially because reducing pump 2 also reduces the indirect supply
through tank 3. This demonstrates the importance of multivariable coupling.
The observed recovery from saturation does not establish robustness to
persistent saturation; an anti-windup design would be needed to address it.

\subsection{Independent linear and nonlinear feedback loops}

The comparison uses $u_s=[300,300]^T$, $d_s=0$, and initial lower-tank
levels $[24.672,36.855]^T\ \mathrm{cm}$. The mass reference for tank 1
increases by $2000\ \mathrm{g}$ at $200\ \mathrm{s}$, followed by the same
increase for tank 2 at $1000\ \mathrm{s}$. Each increase corresponds to
$2000/(\rho A_i)=5.261\ \mathrm{cm}$.
The linear and nonlinear plants have separate PI integrator states and
separate feedback loops. Both plant states use forward Euler with
$T_s=0.5\ \mathrm{s}$ over $2000\ \mathrm{s}$. Each loop computes
\begin{equation}
 I_{k+1}=I_k+T_s e_k,\qquad
 u_k=\operatorname{clip}_{[0,1000]}
 \left(u_s+W_m(K_p e_k+K_i I_{k+1})\right),
\end{equation}
where $K_p,K_i$ are diagonal gain matrices. The linear state is integrated
in deviation coordinates with its own command $u_k-u_s$.

Independent Gaussian sensor errors have standard deviations
$0,10,50\ \mathrm{g}$, equivalent to $0,0.0263,0.1315\ \mathrm{cm}$.
The experiments reuse the same standardized noise array, generated with
seed 100, scaled to each noise level. Linear and nonlinear loops receive
the same noise realization within each case. No process noise is included.
Figure~\ref{fig:p3-comparison} plots true masses, rather than noisy readings.

\begin{table}[htbp]
 \centering
 \caption{Maximum absolute difference between true nonlinear and linear
 levels over the full comparison, and nonlinear pump standard deviations
 over $1500$--$2000\ \mathrm{s}$.}
 \label{tab:p3-comparison}
 \begin{tabular}{rrrrr}
 \hline
 Sensor SD (g) & $\max|h_1^{NL}-h_1^L|$ (cm) &
 $\max|h_2^{NL}-h_2^L|$ (cm) & SD($F_1$) & SD($F_2$)\\
 \hline
 0 & 0.0732 & 0.1169 & 0.0574 & 0.0594\\
 10 & 0.0732 & 0.1169 & 0.3773 & 0.3214\\
 50 & 0.0731 & 0.1169 & 1.8499 & 1.5901\\
 \hline
 \end{tabular}
\end{table}
Pump standard deviations in Table~\ref{tab:p3-comparison} are in
$\mathrm{cm^3/s}$. The noiseless values include the remaining settling drift.
The two models predict similar level trajectories for these steps, while
larger sensor noise increases pump variability. Without noise, the final
nonlinear mass errors are $[0.0668,0.1240]^T\ \mathrm{g}$, corresponding to
less than $0.0004\ \mathrm{cm}$. Agreement in these experiments supports
local use of the linear design, but does not establish performance for
arbitrary operating points, process disturbances, or larger reference changes.

\begingroup
\newcommand{\pThreeFigure}[1]{%
 \IfFileExists{problem03/results/#1.pdf}{%
 \includegraphics[width=\linewidth]{problem03/results/#1.pdf}}{%
 \includegraphics[width=\linewidth]{results/#1.pdf}}}
\begin{figure}[htbp]
 \centering\pThreeFigure{analysis_01}
 \caption{Linear mass responses to unit pump-flow steps at the analysis
 operating point. The direct and cross paths have different dynamics.}
 \label{fig:p3-open}
\end{figure}
\begin{figure}[htbp]
 \centering\pThreeFigure{analysis_02}
 \caption{Full linear mass-reference response with static decoupling and
 PI control. Off-diagonal panels show residual transient interaction.}
 \label{fig:p3-linear}
\end{figure}
\begin{figure}[htbp]
 \centering\pThreeFigure{level_control_01}
 \caption{Linear level-reference responses for the PI design used in the
 nonlinear level experiment.}
 \label{fig:p3-level-linear}
\end{figure}
\begin{figure}[htbp]
 \centering\pThreeFigure{level_control_02}
 \caption{Nonlinear level regulation to $[45,30]^T\ \mathrm{cm}$ using
 sampled PI control. The lower panel shows the actual pump commands.}
 \label{fig:p3-level}
\end{figure}
\begin{figure}[htbp]
 \centering\pThreeFigure{comparison_01}
 \caption{Independent nonlinear and linear mass-feedback loops at three
 measurement-noise levels, with identical noise supplied to both plants
 in each case. Dotted curves show the mass references.}
 \label{fig:p3-comparison}
\end{figure}
\endgroup

\section{Scope of the Controller Evaluation}

The available code in this repository and the source scripts in
\texttt{../MPC/four-tank-model/} implement PI control. P and filtered PID
structures are given above, but their tuning and closed-loop simulations
have not yet been implemented. Consequently, the figures and numerical
results in this chapter evaluate PI control only. Completing the requested
P/PI/PID comparison requires separate P and PID designs tested on the same
references, operating point, disturbance sequences, sensor noise, and pump
constraints. Such tests should compare tracking error, settling, interaction,
and actuator variability, with derivative filtering and saturation treatment
specified explicitly.
